%package list
\documentclass{article}
\usepackage[top=3cm, bottom=3cm, outer=3cm, inner=3cm]{geometry}
\usepackage{multicol}
\usepackage{graphicx}
\usepackage{url}
%\usepackage{cite}
\usepackage{hyperref}
\usepackage{array}
%\usepackage{multicol}
\newcolumntype{x}[1]{>{\centering\arraybackslash\hspace{0pt}}p{#1}}
\usepackage{natbib}
\usepackage{pdfpages}
\usepackage{multirow}
\usepackage[normalem]{ulem}
\useunder{\uline}{\ul}{}
\usepackage{svg}
\usepackage{xcolor}
\usepackage{listings}
\lstdefinestyle{ascii-tree}{
    literate={├}{|}1 {─}{--}1 {└}{+}1 
  }
\lstset{basicstyle=\ttfamily,
  showstringspaces=false,
  commentstyle=\color{red},
  keywordstyle=\color{blue}
}
%\usepackage{booktabs}
\usepackage{caption}
\usepackage{subcaption}
\usepackage{float}
\usepackage{array}

\newcolumntype{M}[1]{>{\centering\arraybackslash}m{#1}}
\newcolumntype{N}{@{}m{0pt}@{}}


%%%%%%%%%%%%%%%%%%%%%%%%%%%%%%%%%%%%%%%%%%%%%%%%%%%%%%%%%%%%%%%%%%%%%%%%%%%%
%%%%%%%%%%%%%%%%%%%%%%%%%%%%%%%%%%%%%%%%%%%%%%%%%%%%%%%%%%%%%%%%%%%%%%%%%%%%
\newcommand{\itemEmail}{fayamamani@unsa.edu.pe}
\newcommand{\itemStudent}{Francis Willian Ayamamani Añasco}
\newcommand{\itemCourse}{Programación Web 2}
\newcommand{\itemCourseCode}{GUIA-PRLD-001}
\newcommand{\itemSemester}{I}
\newcommand{\itemUniversity}{Universidad Nacional de San Agustín de Arequipa}
\newcommand{\itemFaculty}{Facultad de Ingeniería de Producción y Servicios}
\newcommand{\itemDepartment}{Departamento Académico de Ingeniería de Sistemas e Informática}
\newcommand{\itemSchool}{Escuela Profesional de Ingeniería de Sistemas}
\newcommand{\itemAcademic}{2023 - A}
\newcommand{\itemInput}{Del 5 Junio 2023}
\newcommand{\itemOutput}{Al 5 Junio 2023}
\newcommand{\itemPracticeNumber}{03}
\newcommand{\itemTheme}{Ajax y NodeJS}
%%%%%%%%%%%%%%%%%%%%%%%%%%%%%%%%%%%%%%%%%%%%%%%%%%%%%%%%%%%%%%%%%%%%%%%%%%%%
%%%%%%%%%%%%%%%%%%%%%%%%%%%%%%%%%%%%%%%%%%%%%%%%%%%%%%%%%%%%%%%%%%%%%%%%%%%%

\usepackage[english,spanish]{babel}
\usepackage[utf8]{inputenc}
\AtBeginDocument{\selectlanguage{spanish}}
\renewcommand{\figurename}{Figura}
\renewcommand{\refname}{Referencias}
\renewcommand{\tablename}{Tabla} %esto no funciona cuando se usa babel
\AtBeginDocument{%
	\renewcommand\tablename{Tabla}
}

\usepackage{fancyhdr}
\pagestyle{fancy}
\fancyhf{}
\setlength{\headheight}{30pt}
\renewcommand{\headrulewidth}{1pt}
\renewcommand{\footrulewidth}{1pt}
\fancyhead[L]{\raisebox{-0.2\height}{\includegraphics[width=3cm]{img/logo_episunsa.png}}}
\fancyhead[C]{\fontsize{7}{7}\selectfont	\itemUniversity \\ \itemFaculty \\ \itemDepartment \\ \itemSchool \\ \textbf{\itemCourse}}
\fancyhead[R]{\raisebox{-0.2\height}{\includegraphics[width=1.2cm]{img/logo_abet}}}
\fancyfoot[L]{Estudiante Francis Ayamamani Añasco}
\fancyfoot[C]{\itemCourse}
\fancyfoot[R]{Página \thepage}

% para el codigo fuente
\usepackage{listings}
\usepackage{color, colortbl}
\definecolor{dkgreen}{rgb}{0,0.6,0}
\definecolor{gray}{rgb}{0.5,0.5,0.5}
\definecolor{mauve}{rgb}{0.58,0,0.82}
\definecolor{codebackground}{rgb}{0.95, 0.95, 0.92}
\definecolor{tablebackground}{rgb}{0.8, 0, 0}

\lstset{frame=tb,
	language=bash,
	aboveskip=3mm,
	belowskip=3mm,
	showstringspaces=false,
	columns=flexible,
	basicstyle={\small\ttfamily},
	numbers=none,
	numberstyle=\tiny\color{gray},
	keywordstyle=\color{blue},
	commentstyle=\color{dkgreen},
	stringstyle=\color{mauve},
	breaklines=true,
	breakatwhitespace=true,
	tabsize=3,
	backgroundcolor= \color{codebackground},
}

\begin{document}
	
	\vspace*{10px}
	
	\begin{center}	
		\fontsize{17}{17} \textbf{ Informe de Laboratorio \itemPracticeNumber}
	\end{center}
	\centerline{\textbf{\Large Tema: \itemTheme}}
	%\vspace*{0.5cm}	

	\begin{flushright}
		\begin{tabular}{|M{2.5cm}|N|}
			\hline 
			\rowcolor{tablebackground}
			\color{white} \textbf{Nota}  \\
			\hline 
			     \\[30pt]
			\hline 			
		\end{tabular}
	\end{flushright}	

	\begin{table}[H]
		\begin{tabular}{|x{4.7cm}|x{4.8cm}|x{4.8cm}|}
			\hline 
			\rowcolor{tablebackground}
			\color{white} \textbf{Estudiante} & \color{white}\textbf{Escuela}  & \color{white}\textbf{Asignatura}   \\
			\hline 
			{\itemStudent \par \itemEmail} & \itemSchool & {\itemCourse \par Semestre: \itemSemester \par Código: \itemCourseCode}     \\
			\hline 			
		\end{tabular}
	\end{table}		
	
	\begin{table}[H]
		\begin{tabular}{|x{4.7cm}|x{4.8cm}|x{4.8cm}|}
			\hline 
			\rowcolor{tablebackground}
			\color{white}\textbf{Laboratorio} & \color{white}\textbf{Tema}  & \color{white}\textbf{Duración}   \\
			\hline 
			\itemPracticeNumber & \itemTheme & 04 horas   \\
			\hline 
		\end{tabular}
	\end{table}
	
	\begin{table}[H]
		\begin{tabular}{|x{4.7cm}|x{4.8cm}|x{4.8cm}|}
			\hline 
			\rowcolor{tablebackground}
			\color{white}\textbf{Semestre académico} & \color{white}\textbf{Fecha de inicio}  & \color{white}\textbf{Fecha de entrega}   \\
			\hline 
			\itemAcademic & \itemInput &  \itemOutput  \\
			\hline 
		\end{tabular}
	\end{table}
	
	\section{Tarea}
	\begin{itemize}		
		\item Cree una aplicación NodeJS con express, para administrar una agenda personal.
        \item Home (“/”) : Página Principal
        \item Preferiblemente trabaje todo en una misma interfaz. 
        \item Ejemplo de estructura de la agenda cuando se quiere “Ver eventos”.
agenda[DIR]
        \begin{itemize}
        \begin{lstlisting}[style=ascii-tree]
|------------- 2023-05-15 [DIR]
|		|-------------- 10-00.txt [FILE]
|		|-------------- 13-30.txt [FILE]
|------------- 2023-05-20 [DIR]
	|-------------- 09-10.txt [FILE]
	|-------------- 20-40.txt [FILE]
\end{lstlisting}
        \item \textbf{La aplicación debe permitir:}
		\begin{itemize}
		    \item Crear evento: fecha y hora. (Si ya existe el archivo no debería ingresar el evento)(La primera línea es el título del evento, las demás líneas son la descripción del evento.
            \item Editar evento. (Se muestran el archivo  donde esta el detalle del evento)
            \item Eliminar evento.
            \item Ver eventos. Utilizar el formato árbol especificado anteriormente, donde debería incluirse sólo el título del evento.
		\end{itemize}
		\item \textbf{Recomendaciones:}
        \item Interfaz que trabaja todo un CRUD en una sóla vista. (Se pueden usar ventanas emergentes)
        \end{itemize}
	\section{MARCO CONCEPTUAL}
	\begin{itemize}
		\item Editor Vim
        \begin{itemize}
            \item Vim es un editor de texto muy eficiente para programar en consola que no tiene nada que envidiar a los editores gráficos.
        \end{itemize}
        \item Instalar Vim en
        \begin{itemize}
            \item GNU/Linux
            \begin{itemize}
                \begin{lstlisting}
                    apt-get install vim
                \end{lstlisting}
            \end{itemize}
            \item MS Windows
            \begin{itemize}
                \item Para descarga en sistemas MS Windows \url{https://www.vim.org/download.php}
            \end{itemize}
            \item MacOS
            \begin{itemize}
                \begin{lstlisting}
                    brew install macvim
                \end{lstlisting}
            \end{itemize}
        \end{itemize}
        \item Git
        \begin{itemize}
            \item Git es un sistema de control de versiones creado por Linus Torvalds. Es eficiente, confiable.
        \end{itemize}
        \item Instalar Git
        \begin{itemize}
            \item GNU/Linux
            \begin{itemize}
                \begin{lstlisting}
                apt-get install git    
                \end{lstlisting}
            \end{itemize}
            \item MS Windows
            \begin{itemize}
                \item Para descarga en sistemas MS Windows \url{https://git-scm.com/download/win}
            \end{itemize}
            \item MacOS
            \begin{itemize}
                \item Para descarga en sistemas MacOS \url{https://git-scm.com/download/mac}
                \begin{verbatim}
                    brew install git
                \end{verbatim}
            \end{itemize}
        \end{itemize}
	\end{itemize}
	
	\section{URL de Repositorio Github}
	\begin{itemize}
		\item URL del Repositorio GitHub para clonar o recuperar.
		\item \url{https://github.com/rescobedoq/programacion.git}
		\item URL para el laboratorio 01 en el Repositorio GitHub.
		\item \url{https://github.com/rescobedoq/programacion/tree/main/lab01}
	\end{itemize}
	
	\section{Resolucion de la tarea}
	
	\subsection{Explicacion del Codigo}
    \begin{figure}[H]
		\centering
		\includegraphics[width=0.8\textwidth,keepaspectratio]{img/indexjs1.png}
		%\includesvg{img/automata.svg}
		%\label{img:mot2}
		%\caption{Product backlog.}
	\end{figure}
	En la primera linea hace referencia al uso del paquete express:
    \begin{lstlisting}[language=bash,caption={Instalacion de Express}][H]
		$ npm i express
	\end{lstlisting}
	Luego al paquete de fs, que viene por defecto; pero todos estos, representados por una constante (fs, express). Asi tambien llamamos a la funcion del paquete express en la variable app, la constante port señala el puerto al que va a escuchar el servidor.
    \begin{lstlisting}[language=bash,caption={app.use}][H]
		app.use(express.static('public'));
        app.use(express.urlencoded({extended: true}));
	\end{lstlisting}
    la contante app, en la primera linea establece ala carpeta public como recurso para su ejecucion.Y la siguiente estable el enlace POST.
    \begin{figure}[H]
		\centering
		\includegraphics[width=0.8\textwidth,keepaspectratio]{img/indexjs2.png}
		%\includesvg{img/automata.svg}
		%\label{img:mot2}
		%\caption{Product backlog.}
	\end{figure}
    En las lineas 10 a la 12 se hace llamado a la rama principal que hace imprime el index.html
    En las lineas 15 a la 25 nos refrencia en la direccion '/eventos', para imprimir el listado de elementos en el elemento './eventos' que se encuentra como argumento en la funcion de fs.readdir()
    \begin{figure}[H]
		\centering
		\includegraphics[width=0.8\textwidth,keepaspectratio]{img/indexjs3.png}
		%\includesvg{img/automata.svg}
		%\label{img:mot2}
		%\caption{Product backlog.}
	\end{figure}
    En las lineas 41 a la 43 nos refrencia al metodo POST, crear un archivo en el listado de elementos, definimos una variable como titulo, que es el argumento que nos envian; de igual manera la descripcion.
    (45-61) se ubica en la carpeta de './eventos' para luego entrar en el archivo 'titulo', y verificar su existencia:
    \begin{itemize}
        \item \textbf{TRUE}, le envia el mensaje: "El evento ya existe"
        \item \textbf{FALSE}, concatena el titulo y la descripcion con un salto de linea
        \begin{itemize}
            \item ejecuta la funcion de fs.write, para crear el archivo, agregarle el contenido (Titulo-descripcion)
            \begin{itemize}
                \item \textbf{TRUE}, lo dirige a la rama principal
                \item \textbf{FALSE}, imprime "Error al crear el evento"
            \end{itemize}
        \end{itemize}
    \end{itemize}
    \begin{figure}[H]
		\centering
		\includegraphics[width=0.8\textwidth,keepaspectratio]{img/indexjs4.png}
		%\includesvg{img/automata.svg}
		%\label{img:mot2}
		%\caption{Product backlog.}
	\end{figure}
    En las lineas 4de app.delete remueve la |ruta de ./eventos/:eventos; siendo ':evento' la variable del archivo, y para ello la definimos esta variable dentro de la funcion, con la constante evento.
    \\Luego la funcion fs.unlink remueve el archivo en la direccion "./eventos/evento.txt"
    \begin{itemize}
        \item \textbf{TRUE}, ejecuta la linea que imprimr "Evento eliminado exitosamente"
        \item \textbf{FALSE}, imprime "Error al eliminar el evento"
    \end{itemize}
    Finalmente la linea conocida para determinar el puerto del servidor y su ejecucion
    \subsection{./public/script.js}
    \begin{figure}[H]
		\centering
		\includegraphics[width=0.8\textwidth,keepaspectratio]{img/script1.png}
		%\includesvg{img/automata.svg}
		%\label{img:mot2}
		%\caption{Product backlog.}
	\end{figure}
 \begin{figure}[H]
		\centering
		\includegraphics[width=0.8\textwidth,keepaspectratio]{img/script2.png}
		%\includesvg{img/automata.svg}
		%\label{img:mot2}
		%\caption{Product backlog.}
	\end{figure}
    En este archivo declaro las funciones del archivo html para direccionar a una nueva ruta, y de la misma manera vincular las elementos impregnados de una funcion POST
	\subsection{Estructura de laboratorio 03}
	\begin{itemize}	
		\item El contenido que se entrega en este laboratorio es el siguiente:
	\end{itemize}
	
\begin{lstlisting}[style=ascii-tree]
lab3/
|--- lab03/
    |   |--- eventos
    |   |--- miestilo.css
    |   |--- LATEX
            |--- img
            |   |--- Ejecucion.png
            |   |--- logo_abet.png
            |   |--- index.png
            |   |--- inicio.png
            |   |--- estilo.png
            |   |--- logo_episunsa.png
            |   |--- logo_unsa.jpg
            |   |--- pseudocodigo_insercion.png    
            |--- lab1.pdf    
            |--- lab1.tex
\end{lstlisting}    

\section{Cuestionario}
	\begin{itemize}
		\item En el Ejemplo "Hola Mundo" con NodeJS. ¿Qué pasó con la línea: "Content type ….."?
        \\ No es necesario debido a que la ejecucion solo imprime un mensaje en la pagina principal, y no se redirige 
        \item En los ejercicios. ¿En qué lugar debería estar el archivo poema.txt?
        \\ Respeto a la propiedad intelectual,Colaboración transparente,Respeto a las políticas del proyecto,Reconocimiento y atribución. Al promover estas conductas éticas, se fomenta un entorno de trabajo colaborativo, respetuoso y responsable en el uso de un Sistema de Control de Versiones.
        \item ¿Entiende la expresión regular en el código y se da cuenta de para qué es útil?
        \\Los estándares de codificación son conjuntos de pautas y reglas establecidas para garantizar la consistencia y legibilidad del código fuente en un proyecto de desarrollo de software. Estos estándares definen un conjunto común de convenciones y prácticas que los desarrolladores deben seguir al escribir código.
        \item Note que la respuesta del servidor está en formato JSON, ¿Habrá alguna forma de verla directamente?
        
	\end{itemize}		

	\section{\textcolor{red}{Rúbricas}}
	
	\subsection{\textcolor{red}{Entregable Informe}}
	\begin{table}[H]
		\caption{Tipo de Informe}
		\setlength{\tabcolsep}{0.5em} % for the horizontal padding
		{\renewcommand{\arraystretch}{1.5}% for the vertical padding
		\begin{tabular}{|p{3cm}|p{12cm}|}
			\hline
			\multicolumn{2}{|c|}{\textbf{\textcolor{red}{Informe}}}  \\
			\hline 
			\textbf{\textcolor{red}{Latex}} & \textcolor{blue}{El informe está en formato PDF desde Latex,  con un formato limpio (buena presentación) y facil de leer.}   \\ 
			\hline 
			
			
		\end{tabular}
	}
	\end{table}
	
	\clearpage
	
	\subsection{\textcolor{red}{Rúbrica para el contenido del Informe y demostración}}
	\begin{itemize}			
		\item El alumno debe marcar o dejar en blanco en celdas de la columna \textbf{Checklist} si cumplio con el ítem correspondiente.
		\item Si un alumno supera la fecha de entrega,  su calificación será sobre la nota mínima aprobada, siempre y cuando cumpla con todos lo items.
		\item El alumno debe autocalificarse en la columna \textbf{Estudiante} de acuerdo a la siguiente tabla:
	
		\begin{table}[ht]
			\caption{Niveles de desempeño}
			\begin{center}
			\begin{tabular}{ccccc}
    			\hline
    			 & \multicolumn{4}{c}{Nivel}\\
    			\cline{1-5}
    			\textbf{Puntos} & Insatisfactorio 25\%& En Proceso 50\% & Satisfactorio 75\% & Sobresaliente 100\%\\
    			\textbf{2.0}&0.5&1.0&1.5&2.0\\
    			\textbf{4.0}&1.0&2.0&3.0&4.0\\
    		\hline
			\end{tabular}
		\end{center}
	\end{table}	
	
	\end{itemize}
	
	\begin{table}[H]
		\caption{Rúbrica para contenido del Informe y demostración}
		\setlength{\tabcolsep}{0.5em} % for the horizontal padding
		{\renewcommand{\arraystretch}{1.5}% for the vertical padding
		%\begin{center}
		\begin{tabular}{|p{2.7cm}|p{7cm}|x{1.3cm}|p{1.2cm}|p{1.5cm}|p{1.1cm}|}
			\hline
    		\multicolumn{2}{|c|}{Contenido y demostración} & Puntos & Checklist & Estudiante & Profesor\\
			\hline
			\textbf{1. GitHub} & Hay enlace URL activo del directorio para el  laboratorio hacia su repositorio GitHub con código fuente terminado y fácil de revisar. &2 &X &2 & \\ 
			\hline
			\textbf{2. Commits} &  Hay capturas de pantalla de los commits más importantes con sus explicaciones detalladas. (El profesor puede preguntar para refrendar calificación). &4 &X &2 & \\ 
			\hline 
			\textbf{3. Código fuente} &  Hay porciones de código fuente importantes con numeración y explicaciones detalladas de sus funciones. &2 &X &1 & \\ 
			\hline 
			\textbf{4. Ejecución} & Se incluyen ejecuciones/pruebas del código fuente  explicadas gradualmente. &2 &X &1 & \\ 
			\hline			
			\textbf{5. Pregunta} & Se responde con completitud a la pregunta formulada en la tarea.  (El profesor puede preguntar para refrendar calificación).  &2 &X &2 & \\ 
			\hline	
			\textbf{6. Fechas} & Las fechas de modificación del código fuente estan dentro de los plazos de fecha de entrega establecidos. &2 &X &1 & \\ 
			\hline 
			\textbf{7. Ortografía} & El documento no muestra errores ortográficos. &2 &X &2 & \\ 
			\hline 
			\textbf{8. Madurez} & El Informe muestra de manera general una evolución de la madurez del código fuente,  explicaciones puntuales pero precisas y un acabado impecable.   (El profesor puede preguntar para refrendar calificación).  &4 &X &1 & \\ 
			\hline
			\multicolumn{2}{|c|}{\textbf{Total}} &20 & &12 & \\ 
			\hline
		\end{tabular}
		%\end{center}
		%\label{tab:multicol}
		}
	\end{table}
	
\clearpage

\section{Referencias}
\begin{itemize}			
	\item \url{https://www.w3schools.com/java/default.asp}
	\item \url{https://www.geeksforgeeks.org/insertion-sort/}
    \item \url{https://github.com/WillianAyamamani/pw2-lab-d-23a/tree/main/lab1}
\end{itemize}	
	
%\clearpage
%\bibliographystyle{apalike}
%\bibliographystyle{IEEEtranN}
%\bibliography{bibliography}
			
\end{document}
